\documentclass[11pt]{article}

\usepackage[T1]{fontenc}
\usepackage{amsmath,amssymb,amsthm}
\usepackage{hyperref}

\title{Structured versus Saturated Information in Non-normal FIML MLR}
\author{}
\date{\today}

\newcommand{\E}{\mathbb{E}}
\newcommand{\tr}{\operatorname{tr}}
\newcommand{\avar}{\operatorname{avar}}
\newcommand{\vech}{\operatorname{vech}}
\newcommand{\argmin}{\operatorname*{arg\,min}}
\newcommand{\eig}{\operatorname{eig}}

\newtheorem*{claim}{Claim}
\newtheorem*{conclusion}{Conclusion}

\begin{document}
\maketitle

\section*{Question}

FIML MLR uses an observed-data normal likelihood for rows with arbitrary
missingness patterns. Lavaan exposes an H1-information choice usually described
as structured versus unstructured. The practical question is whether the
structured choice is mathematically correct for non-normal data, or only a
default inherited from correctly specified normal-theory ML.

This note separates three statements:
\begin{enumerate}
  \item the trace shortcut is exact;
  \item structured information is justified under correctly specified normal
        FIML;
  \item non-normal FIML needs extra equalities for structured information.
\end{enumerate}

\section*{Setup}

Let row \(i\) have observed index set \(o_i\). The working FIML model is the
observed-subvector normal log likelihood
\[
  \ell_i(\theta)
    = \log \phi_{|o_i|}
      (x_{i,o_i}; \mu_{o_i}(\theta), \Sigma_{o_i o_i}(\theta)),
\]
with missingness treated as ignorable. The estimator solves the working-score
equation
\[
  \sum_{i=1}^N s_i(\hat\theta)=0,
  \qquad
  s_i(\theta)=\frac{\partial\ell_i(\theta)}{\partial\theta}.
\]

Under ordinary misspecified-likelihood regularity,
\[
  \sqrt N(\hat\theta-\theta^\star)
    \Rightarrow
    N(0, A_\theta^{-1} B_\theta A_\theta^{-1}),
\]
where
\[
  A_\theta = -\E\left[
    \frac{\partial s_i(\theta^\star)}{\partial\theta'}
  \right],
  \qquad
  B_\theta = \E[s_i(\theta^\star)s_i(\theta^\star)'].
\]
If the observed-data normal model is correctly specified, then information
equality gives \(A_\theta=B_\theta\). Under non-normality, generally
\(A_\theta\ne B_\theta\).

\section*{Saturated H1 coordinates}

For FIML goodness-of-fit, introduce a saturated normal working model with
coordinates
\[
  \eta = (\mu, \vech(\Sigma)).
\]
For each observed pattern, the row likelihood uses the matching observed
subvector of \(\eta\). Let
\[
  h_i(\eta)=\frac{\partial\ell_i(\eta)}{\partial\eta}
\]
be the saturated score. The saturated quasi-MLE \(\hat\eta\) converges to
\[
  \eta^\star = \argmin_\eta
    \E[-\ell_i(\eta)],
\]
not necessarily to a parameter of the true distribution, because a finite
normal saturated mean/covariance model is not saturated for arbitrary
non-normal observed data.

Its robust covariance is
\[
  \avar\{\sqrt N(\hat\eta-\eta^\star)\}
    = H_\eta^{-1} J_\eta H_\eta^{-1},
\]
with
\[
  H_\eta = -\E\left[
    \frac{\partial h_i(\eta^\star)}{\partial\eta'}
  \right],
  \qquad
  J_\eta = \E[h_i(\eta^\star)h_i(\eta^\star)'].
\]
This is the missing-data analogue of the complete-data ADF covariance of
sample moments. In magmaan this is the FIML saturated
\(\Gamma_{\mathrm{mis}} = H_\eta^{-1}J_\eta H_\eta^{-1}\) path.

\section*{The exact trace identity}

Let
\[
  \Delta = \frac{\partial \eta(\theta)}{\partial\theta'}
\]
be the derivative of the restricted model's H1 coordinates with respect to the
restricted parameters, evaluated at \(\theta^\star\) or \(\hat\theta\). Let
\(V\) be the Hessian/weight used by the local quadratic test in \(\eta\)-space.
The usual residual-space matrix is
\[
  U = V - V\Delta(\Delta'V\Delta)^{-1}\Delta'V.
\]
For any conforming covariance matrix \(\Gamma\),
\[
  \tr(U\Gamma)
  =
  \tr(V\Gamma)
  -
  \tr\{(\Delta'V\Delta)^{-1}\Delta'V\Gamma V\Delta\}.
\]
This is only cyclic invariance of the trace. It is an identity, not an
approximation. Therefore a structured-versus-unstructured discrepancy cannot
be blamed on avoiding explicit formation of \(U\).

\section*{Correct normal FIML}

\begin{claim}
If the observed-data normal FIML model is correctly specified and the
restricted SEM is true, structured and saturated H1 information agree
asymptotically.
\end{claim}

\noindent
Under correct normality, information equality holds in the saturated H1 model:
\[
  H_\eta = J_\eta.
\]
The saturated covariance is therefore
\[
  \Gamma_{\mathrm{mis}}
    = H_\eta^{-1}J_\eta H_\eta^{-1}
    = H_\eta^{-1}.
\]
The inverse saturated covariance is \(H_\eta\), the normal-theory observed-data
information. If the restricted model is true, then the restricted
model-implied covariance and the saturated covariance have the same limit.
Consequently a structured evaluation of the normal-theory H1 information and
an unstructured/saturated evaluation target the same limiting matrix:
\[
  V_{\mathrm{struct}}
    \to H_\eta
    \leftarrow
  V_{\mathrm{sat}}.
\]

Thus structured information is mathematically justified in the classical
case. This is the case in which the default is cleanest: it is both
asymptotically correct and computationally convenient.

\section*{Non-normal FIML}

Under non-normality, the same observed-data normal likelihood defines a
quasi-MLE. The limiting restricted parameter is a pseudo-true value
\[
  \theta^\star = \argmin_\theta
    \E[-\ell_i\{\mu(\theta),\Sigma(\theta)\}],
\]
and the saturated normal pseudo-true value is
\[
  \eta^\star = \argmin_\eta \E[-\ell_i(\eta)].
\]
There is no general reason for
\[
  \eta(\theta^\star) = \eta^\star.
\]
Even if first and second observed moments are matched in some simple
complete-data cases, missingness patterns make the saturated pseudo-target a
pattern-weighted observed-data object, while the restricted target is forced
through one common SEM covariance and mean structure.

More importantly, information equality fails:
\[
  H_\eta \ne J_\eta.
\]
The covariance of the saturated normal estimator is the sandwich
\[
  H_\eta^{-1}J_\eta H_\eta^{-1},
\]
not \(H_\eta^{-1}\). A trace correction using \(V=H_\eta\) and
\(\Gamma=H_\eta^{-1}J_\eta H_\eta^{-1}\) is a valid quasi-likelihood
weighted-mixture construction. But a structured choice that replaces
\(H_\eta\) by the model-implied normal information at \(\eta(\theta^\star)\)
requires the additional equality
\[
  H_\eta\{\eta^\star\}
  =
  H_\eta\{\eta(\theta^\star)\},
\]
or at least equality after projection by \(U\). This is not guaranteed by
robust M-estimation theory.

\section*{What structured assumes}

In non-normal FIML, the structured H1 convention effectively assumes that the
local saturated curvature may be evaluated at the restricted model-implied
normal covariance:
\[
  V_{\mathrm{struct}}
    = H_\eta\{\eta(\hat\theta)\}.
\]
The saturated/unstructured convention evaluates the H1 curvature at the
saturated pseudo-target:
\[
  V_{\mathrm{sat}}
    = H_\eta(\hat\eta).
\]

The first is a model-based curvature. The second is a quasi-likelihood
saturated curvature. Under correct normality they coincide asymptotically.
Under non-normality, \(J_\eta\) generally differs from \(H_\eta\). The two
curvatures differ asymptotically only when the restricted and saturated
observed-data targets differ, for example through restricted-model
misspecification or through a non-normal MAR quasi-likelihood target.

\begin{conclusion}
Structured H1 information is not an unconditional consequence of robust
non-normal FIML theory. It is justified when the pseudo-true saturated
curvature and the model-implied structured curvature collapse to the same
limit, as under the correctly specified complete-data or MCAR null. Robust
FIML theory by itself supplies the sandwich \(H^{-1}JH^{-1}\), not the
substitution \(H_\eta\{\eta(\theta^\star)\}\) for \(H_\eta(\eta^\star)\).
\end{conclusion}

\section*{Structured MLR analogue versus saturated FMG}

The apples-to-apples comparison is an \(\eta\)-space comparison. Hold fixed
the model derivative
\[
  \Delta = \frac{\partial \eta(\theta)}{\partial\theta'}
\]
and hold fixed the saturated sandwich covariance
\[
  \Gamma_{\mathrm{mis}}
    =
  H_\eta(\hat\eta)^{-1}J_\eta(\hat\eta)H_\eta(\hat\eta)^{-1}.
\]
Then the only difference between the structured H1 analogue and magmaan's
saturated FMG construction is the curvature matrix \(V\) used to form the
residual operator
\[
  U(V) = V - V\Delta(\Delta'V\Delta)^{-1}\Delta'V.
\]

The structured H1 analogue uses
\[
  V_{\mathrm{struct}} = H_\eta\{\eta(\hat\theta)\},
  \qquad
  \lambda_{\mathrm{struct}}
    =
  \eig\{U(V_{\mathrm{struct}})\Gamma_{\mathrm{mis}}\}.
\]
The saturated FMG construction uses
\[
  V_{\mathrm{sat}} = H_\eta(\hat\eta),
  \qquad
  \lambda_{\mathrm{sat}}
    =
  \eig\{U(V_{\mathrm{sat}})\Gamma_{\mathrm{mis}}\}.
\]

Thus, after both methods are put in the same saturated moment space with the
same empirical sandwich \(\Gamma_{\mathrm{mis}}\), the substantive
structured-versus-saturated distinction is exactly the evaluation point for
\(V\). MLR reporting usually collapses this object to a trace or mean scaling,
for example
\[
  c = \frac{1}{\mathrm{df}}\tr\{U(V)\Gamma_{\mathrm{mis}}\},
\]
whereas FMG keeps the full weighted-mixture spectrum
\(\lambda_1,\ldots,\lambda_{\mathrm{df}}\). The scalar trace is therefore a
summary of the same \(U\Gamma\) object, not a different algebra.

\section*{Interpretation for the MLR problem}

This gives a precise version of the empirical hypothesis: MLR can fail because
it robustifies the covariance but keeps the wrong H1 curvature.
The meat \(J\) sees the non-normal observed-data scores. The structured H1
weight still uses the restricted model-implied normal curvature. That is a
hybrid: empirical score covariance paired with a curvature from the wrong
observed-data geometry.

The unstructured/saturated route does not make the normal working likelihood
true. It only removes one unnecessary model-implied substitution from the H1
side of the trace. Therefore it can improve calibration without solving every
problem caused by treating non-normal or discretized data as continuous.

\section*{Failure channels}

There are two distinct ways for the structured and saturated curvatures to
target different limits.

The first is restricted-model misspecification:
\[
  \eta(\theta^\star) \ne \eta^\star.
\]
This channel does not require non-normal data. It can occur under a normal
complete-data distribution if the restricted SEM is false. In that case the
structured H1 curvature follows the restricted pseudo-true covariance, while
the saturated H1 curvature follows the unrestricted covariance.

The second is the missing-data quasi-likelihood channel. Under non-normal MAR
FIML, the saturated observed-data normal quasi-MLE need not target the same
full-population moment object as the restricted model, even when the SEM is a
correct description of the full-data first two moments. The saturated H1
target is determined by the observed-pattern normal likelihood, and the
restricted target is forced through one common mean/covariance structure.

To disprove universal correctness of structured H1 information, it is enough
to exhibit a data-generating distribution and missingness mechanism for which
\[
  H_\eta\{\eta^\star\}
  \ne
  H_\eta\{\eta(\theta^\star)\}
\]
and the difference survives the residual projection \(U\). Complete-data
misspecification examples establish the first channel. MAR non-normal examples
are the sharper null-calibration channel, because they can make the saturated
and restricted observed-data targets diverge without relying on the restricted
SEM being false.

Under a correct restricted model with complete data, or with MCAR missingness
that leaves the saturated and restricted moment targets aligned, the
structured and saturated H1 curvatures should agree asymptotically. In that
regime, non-normality still matters through \(J_\eta\), but the
structured-versus-saturated evaluation point is not the asymptotic source of a
gap.

To prove structured correctness for a particular simulation cell, one would
need to show the stronger projected equality
\begin{align*}
  &\tr\!\left[
    U\{H_\eta(\eta^\star)\}
    H_\eta(\eta^\star)^{-1}J_\eta
    H_\eta(\eta^\star)^{-1}
  \right] \\
  &\qquad =
  \tr\!\left[
    U\{H_\eta(\eta(\theta^\star))\}
    H_\eta(\eta^\star)^{-1}J_\eta
    H_\eta(\eta^\star)^{-1}
  \right],
\end{align*}
or an equivalent equality for the full weighted-mixture spectrum. There is no
general theorem giving this under non-normal FIML.

\section*{A two-variable witness}

A maximally small complete-data misspecification example shows the first
failure channel. Let the saturated covariance of two observed variables be
\[
  \Sigma_Y =
  \begin{bmatrix}
    1 & \rho \\
    \rho & 1
  \end{bmatrix},
  \qquad \rho \ne 0,
\]
but fit the restricted independence model
\[
  \Sigma(\theta) =
  \begin{bmatrix}
    \theta_1 & 0 \\
    0 & \theta_2
  \end{bmatrix}.
\]
The pseudo-true restricted covariance is
\[
  \Sigma(\theta^\star) = I_2,
\]
whereas the saturated covariance is \(\Sigma_Y\). Hence the structured and
unstructured H1 weights evaluate the normal-theory covariance-moment weight at
different places.

This witness does not rely on non-normality; it isolates the evaluation-point
problem. The same gap arises for normal data if the restricted independence
model is false. Its role is therefore to show that structured H1 information
is not an algebraic consequence of the trace formula, not to prove that
non-normality alone creates an H1-curvature gap under a correct complete-data
or MCAR null.

Use covariance moments ordered as
\[
  \eta = (\sigma_{11},\sigma_{21},\sigma_{22})'.
\]
The restricted derivative is
\[
  \Delta =
  \begin{bmatrix}
    1 & 0 \\
    0 & 0 \\
    0 & 1
  \end{bmatrix},
\]
so the single degree of freedom is the covariance coordinate. At the structured
point \(I_2\),
\[
  \Gamma_{\mathrm{NT}}(I_2)
    =
  \begin{bmatrix}
    2 & 0 & 0 \\
    0 & 1 & 0 \\
    0 & 0 & 2
  \end{bmatrix},
  \qquad
  V_{\mathrm{struct}} = \Gamma_{\mathrm{NT}}(I_2)^{-1}.
\]
The corresponding residual-space matrix is simply
\[
  U_{\mathrm{struct}}
    =
  \begin{bmatrix}
    0 & 0 & 0 \\
    0 & 1 & 0 \\
    0 & 0 & 0
  \end{bmatrix}.
\]

If the trace meat is the normal covariance-moment matrix at the saturated
covariance,
\[
  \Gamma_{\mathrm{NT}}(\Sigma_Y)
    =
  \begin{bmatrix}
    2 & 2\rho & 2\rho^2 \\
    2\rho & 1+\rho^2 & 2\rho \\
    2\rho^2 & 2\rho & 2
  \end{bmatrix},
\]
then the structured one-df scale is
\[
  c_{\mathrm{struct}}
    =
  \tr\{U_{\mathrm{struct}}\Gamma_{\mathrm{NT}}(\Sigma_Y)\}
    =
  1+\rho^2.
\]

By contrast, the unstructured choice uses
\[
  V_{\mathrm{unstruct}}
    = \Gamma_{\mathrm{NT}}(\Sigma_Y)^{-1}.
\]
For any positive definite \(\Gamma\), the matrix
\[
  U(V=\Gamma^{-1})\Gamma
\]
is the \(\Gamma^{-1}\)-orthogonal residual projector and has trace equal to its
rank. Here the rank is the model df, one. Hence
\[
  c_{\mathrm{unstruct}} = 1.
\]

Therefore, for every \(\rho\ne 0\),
\[
  c_{\mathrm{struct}} = 1+\rho^2
  \qquad \text{but} \qquad
  c_{\mathrm{unstruct}} = 1.
\]
The disagreement is not numerical and not a trace shortcut artifact. It is
caused entirely by evaluating \(V\) at the restricted structured covariance
\(I_2\) instead of at the saturated covariance \(\Sigma_Y\).

This witness uses complete data only to keep the algebra visible. In the
FIML saturated construction, missing data replaces
\(\Gamma_{\mathrm{NT}}(\Sigma_Y)\) by the saturated observed-pattern sandwich
\[
  H_\eta^{-1}J_\eta H_\eta^{-1},
\]
but the same logical burden remains: structured correctness requires the
restricted and saturated H1 curvatures to agree after projection.

\section*{Practical diagnostic}

The clean diagnostic is therefore not a new trace formula. It is a controlled
change in the evaluation point for \(V\):
\[
  V_{\mathrm{struct}} = H_\eta\{\eta(\hat\theta)\},
  \qquad
  V_{\mathrm{sat}} = H_\eta(\hat\eta),
\]
with the same empirical score meat and the same trace identity. If
structured follows the bad MLR rejection behavior while saturated/unstructured
moves toward nominal, the structured H1 information convention is the culprit.
If both fail similarly, the failure is deeper in the normal working likelihood
or in the base FIML LRT, not only in the H1-information convention.

\end{document}
